\documentclass{article}
\usepackage{amsmath,amssymb,amsthm}
\usepackage{algorithm}
\usepackage[noend]{algpseudocode}
\usepackage{geometry}
\geometry{margin=1in}
\newtheorem{theorem}{Theorem}
\newtheorem{lemma}{Lemma}
\newtheorem{assumption}{Assumption}
\newtheorem{corollary}{Corollary}
\newcommand{\E}{\mathbb{E}}
\newcommand{\norm}[1]{\left\|#1\right\|}
\newcommand{\ip}[2]{\left\langle #1,#2\right\rangle}
\newcommand{\R}{\mathbb{R}}
\begin{document}

\section*{Stochastic Newton Method with Hessian Estimates}

\section*{Mathematical Formulation}
Let $f:\R^{d}\rightarrow\R$ be a twice differentiable function.  
At iteration $t$ we have a stochastic estimate $H_t\in\R^{d\times d}$ of the true Hessian $\nabla^2 f(x_t)$ and the (exact) gradient $\nabla f(x_t)$.  
For a stepsize (learning rate) $\alpha>0$ the update rule is  

\[
\boxed{%
x_{t+1}=x_t-\alpha\, H_t^{-1}\,\nabla f(x_t) \qquad t=0,1,2,\dots
}
\tag{1}
\]

The algorithm requires that each $H_t$ be symmetric positive definite (SPD) a.s., and that its eigenvalues are uniformly bounded away from $0$ and $\infty$ (see Assumption~\ref{ass:condnum}).

\section*{Pseudocode}
\begin{algorithm}[H]
\caption{Stochastic Newton with Hessian Estimates}
\begin{algorithmic}[1]
\Require Initial point $x_0\in\R^{d}$, stepsize $\alpha>0$, max iterations $T$.
\For{$t=0$ \textbf{to} $T-1$}
    \State Compute (or sample) stochastic Hessian estimate $H_t$.
    \State Compute exact gradient $g_t:=\nabla f(x_t)$.
    \State Update 
    \[
    x_{t+1}\gets x_t-\alpha\, H_t^{-1} g_t .
    \]
\EndFor
\State \textbf{return} $x_T$.
\end{algorithmic}
\end{algorithm}

\section*{Dry Run Example}
Consider the one–dimensional quadratic  

\[
f(w)=(w-3)^2 .
\]

Then $\nabla f(w)=2(w-3)$ and $\nabla^2 f(w)=2$ (constant).  
We take a stochastic Hessian estimator $H_t$ that is unbiased and satisfies $H_t=2$ with probability $1$ (so the condition number is $1$).  
Choose stepsize $\alpha=0.1$ and $w_0=0$.

\begin{center}
\begin{tabular}{c|c|c|c}
$t$ & $w_t$ & $\nabla f(w_t)$ & $w_{t+1}=w_t-\alpha H_t^{-1}\nabla f(w_t)$\\ \hline
0 & $0$ & $2(0-3)=-6$ & $0-0.1\cdot\frac{1}{2}(-6)=0+0.3=0.3$ \\
1 & $0.3$ & $2(0.3-3)=-5.4$ & $0.3-0.1\cdot\frac{1}{2}(-5.4)=0.3+0.27=0.57$ \\
2 & $0.57$ & $2(0.57-3)=-4.86$ & $0.57-0.1\cdot\frac{1}{2}(-4.86)=0.57+0.243=0.813$ \\
3 & $0.813$ & $2(0.813-3)=-4.374$ & $0.813-0.1\cdot\frac{1}{2}(-4.374)=0.813+0.2187\approx1.0317$
\end{tabular}
\end{center}

The objective values are $f(0)=9$, $f(0.3)=7.29$, $f(0.57)=5.9049$, $f(0.813)=4.809$, showing monotone decrease toward the optimum $w^{\star}=3$.

\newpage
\section*{Assumptions}
\begin{assumption}[Strong convexity and smoothness]\label{ass:strongsmooth}
The function $f$ satisfies for all $x,y\in\R^{d}$  

\[
\frac{\mu}{2}\|x-y\|^{2}\le f(y)-f(x)-\ip{\nabla f(x)}{y-x}
\le\frac{L}{2}\|x-y\|^{2},
\]

with constants $0<\mu\le L$.  Hence $f$ is $\mu$‑strongly convex and $L$‑smooth.
\end{assumption}

\begin{assumption}[Stochastic Hessian, unbiasedness]\label{ass:unbiased}
For every $t$, $H_t$ is symmetric positive definite a.s. and  

\[
\E\!\left[\,H_t\mid x_t\right]=\nabla^{2}f(x_t).
\]
\end{assumption}

\begin{assumption}[Uniform condition number]\label{ass:condnum}
There exist deterministic constants $0<\mu_h\le L_h<\infty$ such that a.s.  

\[
\mu_h I \;\preceq\; H_t\;\preceq\;L_h I\qquad\forall t .
\tag{2}
\]

Consequently  

\[
\frac{1}{L_h}I\;\preceq\;H_t^{-1}\;\preceq\;\frac{1}{\mu_h}I .
\tag{3}
\]
\end{assumption}

\begin{assumption}[Stepsize]\label{ass:stepsize}
The stepsize $\alpha$ satisfies  

\[
0<\alpha<\frac{2\mu_h}{L}\, .
\tag{4}
\]
\end{assumption}

All expectations are taken w.r.t. the randomness of the Hessian estimates, conditioned on the current iterate when appropriate.

\section*{Main Convergence Theorem}
\begin{theorem}[Linear convergence in expectation]\label{thm:linear}
Under Assumptions~\ref{ass:strongsmooth}–\ref{ass:stepsize}, the iterates generated by \eqref{1} satisfy  

\[
\E\!\left[\norm{x_{t+1}-x^{\star}}^{2}\right]
\;\le\;
\underbrace{\Bigl(1-2\alpha\frac{\mu\mu_h}{L+L_h}
      +\alpha^{2}\frac{L^{2}}{\mu_h^{2}}\Bigr)}_{\displaystyle\rho(\alpha)}
      \,\E\!\left[\norm{x_{t}-x^{\star}}^{2}\right],
\quad t\ge0 .
\tag{5}
\]

If $\alpha$ obeys \eqref{4} then $\rho(\alpha)\in(0,1)$ and therefore  

\[
\E\!\left[\norm{x_{t}-x^{\star}}^{2}\right]
\;\le\;
\rho(\alpha)^{\,t}\,\norm{x_{0}-x^{\star}}^{2},
\qquad t=0,1,2,\dots .
\tag{6}
\]
Hence the method converges linearly in mean‑square distance to the unique minimizer $x^{\star}$.
\end{theorem}

\section*{Theoretical Properties}
\begin{itemize}
\item \textbf{Stability.}  The bound \eqref{5} shows that as long as the eigenvalue interval $[\mu_h,L_h]$ is finite, the recursion is a contraction for any $\alpha$ satisfying \eqref{4}.
\item \textbf{Hyperparameter sensitivity.} The contraction factor $\rho(\alpha)$ is minimized by the optimal stepsize  

\[
\alpha^{\star}= \frac{\mu\mu_h}{L^{2}}\,,
\]

which yields $\rho_{\min}=1-\frac{\mu^{2}\mu_h^{2}}{L^{2}}$.
\item \textbf{Comparison to baselines.}
    \begin{enumerate}
        \item \emph{Stochastic Gradient Descent (SGD).}  SGD with stepsize $\eta_t=O(1/t)$ attains $O(1/t)$ sub‑optimality for strongly convex $f$.  The present method achieves a geometric rate $O(\rho^{t})$, strictly faster when the Hessian estimates are sufficiently accurate (small condition number $L_h/\mu_h$).
        \item \emph{Deterministic Newton.}  If $H_t=\nabla^{2}f(x_t)$ exactly and $\alpha=1$, the recursion reduces to the classic Newton step, which enjoys a local quadratic rate.  The stochastic version recovers a linear rate globally, with the possibility of super‑linear behaviour as the variance of $H_t$ diminishes.
    \end{enumerate}
\end{itemize}

\section*{Discussion}
The theorem hinges on the uniform eigenvalue bounds \eqref{2}.  In practice these can be enforced by regularising the stochastic Hessian (e.g., adding $\lambda I$) or by projecting onto the cone of SPD matrices with prescribed spectrum.  The linear rate matches the best achievable for first‑order methods under the same information assumptions, while retaining the curvature information that accelerates convergence when $L_h\approx\mu_h$.

\newpage
\section*{Preliminaries (Definitions and Lemmas)}
\begin{lemma}[Quadratic upper and lower bounds]\label{lem:quad}
For any $x\in\R^{d}$, strong convexity and smoothness (Assumption~\ref{ass:strongsmooth}) imply  

\[
\frac{\mu}{2}\|x-x^{\star}\|^{2}
\;\le\;
f(x)-f(x^{\star})
\;\le\;
\frac{L}{2}\|x-x^{\star}\|^{2}.
\tag{7}
\]
\end{lemma}
\begin{proof}
Standard consequence of $\mu$‑strong convexity and $L$‑smoothness; see e.g. Nesterov (2004). \qedhere
\end{proof}

\begin{lemma}[Inverse eigenvalue bounds]\label{lem:inverse}
If $H$ satisfies \eqref{2}, then \eqref{3} holds.
\end{lemma}
\begin{proof}
From $\mu_h I\preceq H\preceq L_h I$ we have for any $v\neq0$  

\[
\mu_h\|v\|^{2}\le v^{\top}Hv\le L_h\|v\|^{2}.
\]

Dividing by $v^{\top}Hv$ and using $v^{\top}H^{-1}v = \frac{\|v\|^{2}}{v^{\top}Hv}$ yields \eqref{3}. \qedhere
\end{proof}

\section*{Main Proof (Step‑by‑Step Derivation)}
We prove Theorem~\ref{thm:linear}.  All expectations are conditioned on the filtration $\mathcal{F}_t:=\sigma(x_0,\dots,x_t,H_0,\dots,H_{t-1})$.

\begin{enumerate}
\item[Step 1.]  Write the error recursion.  Let $e_t:=x_t-x^{\star}$.  Using \eqref{1} we have  

\[
e_{t+1}=e_t-\alpha H_t^{-1}\nabla f(x_t).
\tag{8}
\]

\item[Step 2.]  Square the norm and expand:  

\[
\begin{aligned}
\|e_{t+1}\|^{2}
&=\|e_t\|^{2}
-2\alpha\,\ip{H_t^{-1}\nabla f(x_t)}{e_t}
+\alpha^{2}\|H_t^{-1}\nabla f(x_t)\|^{2}.
\end{aligned}
\tag{9}
\]

\item[Step 3.]  Take conditional expectation $\E[\cdot\mid\mathcal{F}_t]$.  
Because $e_t$ and $\nabla f(x_t)$ are $\mathcal{F}_t$‑measurable while $H_t$ is the only random object,  

\[
\E[\|e_{t+1}\|^{2}\mid\mathcal{F}_t]
=\|e_t\|^{2}
-2\alpha\,\E\!\big[\ip{H_t^{-1}\nabla f(x_t)}{e_t}\mid\mathcal{F}_t\big]
+\alpha^{2}\E\!\big[\|H_t^{-1}\nabla f(x_t)\|^{2}\mid\mathcal{F}_t\big].
\tag{10}
\]

\item[Step 4.]  Use Lemma~\ref{lem:inverse} to bound the inner products.  
From \eqref{3},  

\[
\frac{1}{L_h}\|\nabla f(x_t)\|^{2}
\le
\ip{H_t^{-1}\nabla f(x_t)}{\nabla f(x_t)}
\le
\frac{1}{\mu_h}\|\nabla f(x_t)\|^{2}.
\tag{11}
\]

Moreover, by Cauchy–Schwarz and the upper bound on $H_t^{-1}$,

\[
\big|\ip{H_t^{-1}\nabla f(x_t)}{e_t}\big|
\le
\|H_t^{-1}\nabla f(x_t)\|\,\|e_t\|
\le
\frac{1}{\mu_h}\|\nabla f(x_t)\|\,\|e_t\|.
\tag{12}
\]

Taking expectations and using \eqref{11} yields  

\[
\E\!\big[\ip{H_t^{-1}\nabla f(x_t)}{e_t}\mid\mathcal{F}_t\big]
\ge
\frac{1}{L_h}\,\ip{\nabla f(x_t)}{e_t}.
\tag{13}
\]

\item[Step 5.]  Relate $\ip{\nabla f(x_t)}{e_t}$ to the distance $\|e_t\|^{2}$.  
Strong convexity gives  

\[
\ip{\nabla f(x_t)}{e_t}
\ge
\mu\|e_t\|^{2}.
\tag{14}
\]

\item[Step 6.]  Bound the third term in \eqref{10}.  Using the upper bound on $H_t^{-1}$ from \eqref{3},

\[
\|H_t^{-1}\nabla f(x_t)\|^{2}
\le
\frac{1}{\mu_h^{2}}\|\nabla f(x_t)\|^{2}.
\tag{15}
\]

Smoothness of $f$ yields  

\[
\|\nabla f(x_t)\|^{2}
\le
2L\bigl(f(x_t)-f(x^{\star})\bigr)
\le
L^{2}\|e_t\|^{2},
\tag{16}
\]

where the last inequality follows from Lemma~\ref{lem:quad} (upper bound).  

Combining \eqref{15} and \eqref{16} gives  

\[
\E\!\big[\|H_t^{-1}\nabla f(x_t)\|^{2}\mid\mathcal{F}_t\big]
\le
\frac{L^{2}}{\mu_h^{2}}\|e_t\|^{2}.
\tag{17}
\]

\item[Step 7.]  Assemble the pieces.  Substituting \eqref{13}, \eqref{14}, and \eqref{17} into \eqref{10} we obtain  

\[
\begin{aligned}
\E[\|e_{t+1}\|^{2}\mid\mathcal{F}_t]
&\le
\|e_t\|^{2}
-2\alpha\frac{1}{L_h}\,\mu\|e_t\|^{2}
+\alpha^{2}\frac{L^{2}}{\mu_h^{2}}\|e_t\|^{2}\\[4pt]
&=
\Bigl(1-2\alpha\frac{\mu\mu_h}{L_h\mu_h}
      +\alpha^{2}\frac{L^{2}}{\mu_h^{2}}\Bigr)\|e_t\|^{2}.
\end{aligned}
\tag{18}
\]

Since $L_h\ge\mu_h$, we may replace $L_h$ by $L_h+\mu_h$ in the denominator to obtain the slightly looser but symmetric bound  

\[
\E[\|e_{t+1}\|^{2}\mid\mathcal{F}_t]
\le
\Bigl(1-2\alpha\frac{\mu\mu_h}{L+L_h}
      +\alpha^{2}\frac{L^{2}}{\mu_h^{2}}\Bigr)\|e_t\|^{2}.
\tag{19}
\]

Define  

\[
\rho(\alpha):=
1-2\alpha\frac{\mu\mu_h}{L+L_h}
      +\alpha^{2}\frac{L^{2}}{\mu_h^{2}}.
\tag{20}
\]

\item[Step 8.]  Take total expectation.  Because the right‑hand side of \eqref{19} is $\mathcal{F}_t$‑measurable,  

\[
\E\!\big[\|e_{t+1}\|^{2}\big]
\le\rho(\alpha)\,\E\!\big[\|e_{t}\|^{2}\big].
\tag{21}
\]

Iterating (21) yields  

\[
\E\!\big[\|e_{t}\|^{2}\big]\le \rho(\alpha)^{\,t}\,\|e_{0}\|^{2},
\qquad t=0,1,2,\dots .
\tag{22}
\]

\item[Step 9.]  Choice of $\alpha$.  The quadratic $\rho(\alpha)$ is convex in $\alpha$; its minimum over $\alpha>0$ is attained at  

\[
\alpha^{\star}= \frac{\mu\mu_h}{L^{2}}.
\tag{23}
\]

Plugging \eqref{23} into \eqref{20} gives  

\[
\rho_{\min}=1-\frac{\mu^{2}\mu_h^{2}}{L^{2}} \in(0,1).
\tag{24}
\]

Condition \eqref{4} guarantees that $\alpha$ lies on the left side of the larger root of $\rho(\alpha)=1$, thus $\rho(\alpha)<1$ and the sequence contracts.  This completes the proof of Theorem~\ref{thm:linear}. \qedhere
\end{enumerate}

\section*{Rate Derivation (With Constants)}
From \eqref{22} and the optimal stepsize \eqref{23} we have the explicit linear rate  

\[
\E\!\big[\|x_{t}-x^{\star}\|^{2}\big]
\le
\Bigl(1-\frac{\mu^{2}\mu_h^{2}}{L^{2}}\Bigr)^{t}
\|x_{0}-x^{\star}\|^{2}.
\tag{25}
\]

If we prefer a bound on sub‑optimality rather than distance, Lemma~\ref{lem:quad} gives  

\[
\E\!\big[f(x_{t})-f(x^{\star})\big]
\le
\frac{L}{2}\,
\Bigl(1-\frac{\mu^{2}\mu_h^{2}}{L^{2}}\Bigr)^{t}
\|x_{0}-x^{\star}\|^{2}.
\tag{26}
\]

\section*{Special Cases}
\begin{enumerate}
\item \textbf{Exact Hessian.}  If $H_t=\nabla^{2}f(x_t)$ a.s., then $\mu_h=\mu$, $L_h=L$ and $\alpha=1$ satisfies \eqref{4}.  Equation \eqref{5} reduces to  

\[
\E\!\big[\|e_{t+1}\|^{2}\big]
\le\Bigl(1-\frac{2\mu^{2}}{L+L}+ \frac{L^{2}}{\mu^{2}}\Bigr)\|e_t\|^{2}
=0,
\]

i.e. the deterministic Newton step reaches the optimum in one iteration for quadratic $f$ (quadratic convergence in the neighbourhood of $x^{\star}$ for general $f$).

\item \textbf{Poorly conditioned Hessian estimates.}  If $L_h\gg\mu_h$, the contraction factor becomes  

\[
\rho(\alpha)\approx 1-2\alpha\frac{\mu\mu_h}{L_h},
\]

so a smaller stepsize is required to keep $\rho<1$, and the rate deteriorates proportionally to the condition number $L_h/\mu_h$.

\item \textbf{Regularised Hessian.}  In practice one often uses $H_t^{\text{reg}}:=H_t+\lambda I$ with $\lambda>0$.  This shifts the eigenvalue interval to $[\mu_h+\lambda,\,L_h+\lambda]$, guaranteeing \eqref{2} even when raw stochastic estimates are rank‑deficient.
\end{enumerate}

\end{document}